Although the geometrical VoF algorithm supports arbitrary cell shapes, the mesh
topology dictates the concrete strategy and conceptual approach (i.e.\ type of
topological operations) for mesh refinement. We are applying local dynamic AMR
to increase the solution accuracy especially in the vicinity of interfaces, in
particular for the coupling between the geometrical algorithm and the flow
solution. The algorithm for local dynamic AMR in OpenFOAM is developed on the
basis of the unstructured \emph{hexahedral} mesh: it is based on a \emph{fully
unstructured mesh refinement} procedure. Fully unstructured mesh refinement has
the advantage to deal with flow domains of arbitrary geometrical complexity,
since the topological operations on the mesh are done directly, changing the
overall mesh connectivity information. However, since the connectivity
information is changed by the topological operation, the mesh data structures
need to be modified in order to store the new connectivity produced by the
refinement cycle.  

If the topological operation introduces a production of mesh data, the capacity
of the data structures needs to be extended. This involves re-allocating memory
and extending the capacity of the data structures, which sets a computational
price to the flexibility in dealing with arbitrary complex geometries. Even so,
the mesh refinement engine in OpenFOAM, together with the mapping of all
intensive properties, takes up only a small fraction of the overall time
involved in the flow solution.

The refinement cycle consists of two sub-cycles which may be executed,
depending on the refinement criteria: refinement and un-refinement sub-cycle.
If both sub-cycles are executed, the mesh information from the beginning of the
refinement cycle is lost in between the sub-cycle execution. Removal of
intermediate information makes it impossible to map the mesh modified by two
subsequent refinement cycles to the initial un-refined mesh to which the
interface geometry maps. The obstructed mapping renders the refinement
procedure unusable together with the geometrical transport algorithm. For this
purpose, we have extended the dynamic mesh capability to store the connectivity
information relative to the unmodified mesh state at the beginning of the
refinement cycle. This further enables the execution of a \emph{geometrical
mapping algorithm} for the volume fraction field described by Algorithm
\ref{alg:geometrical-mapping}. 

\begin{algorithm}[H]
    \centering
    \caption{Geometrical mapping algorithm for the volume fraction field}
    \label{alg:geometrical-mapping}

    {\small
      \begin{algorithmic}
        \State initialize the absolute cell map 
        \If{cells to be refined exist}
           \State mark mesh as refined
            \State refine mesh 
            \State absolute cell map = current cell map
        \EndIf
        \If{cells to be unrefined exist}
           \State mark mesh as unrefined 
           \State copy the absolute cell map 
           \State un-refine mesh 
        \EndIf
        \If {mesh is refined and unrefined}
            \For{current cell map} 
               \State absolute cell map [cell] = absolute cell map[current cell map[cell]]
            \EndFor
        \ElsIf{mesh is unrefined}
            \State absolute cell map = current cell map
        \EndIf
        \For{absolute cell map}
            \State intersection = intersect (interface planes [absolute cell map [cell]], cell) 
            \State volume fraction [cell] = intersection / cell volume 
        \EndFor
      \end{algorithmic}
    }
\end{algorithm}

The interested reader is directed to our forthcoming publication of the authors for details
regarding the description of the unstructured mesh topology in OpenFOAM and its
reflection onto the geometrical mapping algorithm for the volume fraction field
\citep*{MaricVoFoamMapping2013}.
